% Was three build-up frames using the paper's mu(s|theta^mu) / Q(s,a|theta^Q)
% conditioning-bar notation. Now one frame in course notation: theta = actor
% (Days 3-4), phi = critic (Day 4), superscript minus = target (Day 2).
\begin{frame}{DDPG Problem Setting: Four Networks}
\vspace{0.8em}
\begin{center}
\renewcommand{\arraystretch}{1.9}
\begin{tabular}{ll}
$\mu_\theta(s)$        & \begin{tabular}{@{}l@{}}\textbf{Actor} network --- deterministic, outputs an action in $\mathbb{R}^d$\end{tabular} \\
$Q_\phi(s,a)$          & \textbf{Critic} network --- action-value \\
$\mu_{\theta^{-}}(s)$  & \textbf{Target actor} --- a slowly-following copy of $\mu_\theta$ \\
$Q_{\phi^{-}}(s,a)$    & \textbf{Target critic} --- a slowly-following copy of $Q_\phi$ \\
$\mathcal{D}$          & \textbf{Replay buffer} of transitions $(s, a, r, s')$ \\
\end{tabular}
\end{center}
\vspace{1em}
\footnotesize
Target networks appear twice because the TD target $y = r + \gamma\, Q_{\phi^{-}}(s', \mu_{\theta^{-}}(s'))$ uses \emph{both} of them --- DQN only needed one, because it had no actor.
\end{frame}
